% Part: normal-modal-logic
% Chapter: syntax-and-semantics
% Section: normal-modal-logics

\documentclass[../../../include/open-logic-section]{subfiles}

\begin{document}

\olfileid{nml}{syn}{nor}

\olsection{منطق‌های موجهاتی نرمال}

می‌توان منطق موجهات را هر مجموعهٔ به‌نحو مناسبی خوش‌رفتار از
!!{formula}ها دانست. برخی منطق‌های موجهات هم به دلیل تفسیرها و
کاربردهای جالبشان و هم، مهم‌تر از آن، به دلیل فهم نظری خوبی که از
آن‌ها داریم، اهمیت ویژه دارند. این‌ها همان منطق‌های موجهاتی موسوم به
نرمال‌اند.

\begin{defn}\ollabel{defn:modal-logic}
یک \emph{منطق موجهات}~$\Log L$ مجموعه‌ای از !!{formula}های زبان پایهٔ
موجهات است که:
\begin{enumerate}
\item همهٔ همان‌گویی‌های گزاره‌ای را در بر دارد؛
\item تحت جانشینی یکنواخت بسته است؛ و
\item تحت وضع مقدم بسته است؛ یعنی هرگاه $!A \in \Log L$ و
  $(!A \lif !B) \in \Log L$، آنگاه $!B \in \Log L$ نیز.
\end{enumerate}
\end{defn}

\begin{ex}
برخی مجموعه‌های نه‌چندان جالب از !!{formula}ها که این تعریف را برآورده
می‌کنند و بنابراین منطق موجهات به شمار می‌آیند عبارت‌اند از:
\begin{enumerate}
\item مجموعهٔ همهٔ همان‌گویی‌ها؛
\item مجموعهٔ همهٔ !!{formula}ها، یعنی منطق موجهات ناسازگار.
\end{enumerate}
\end{ex}

\begin{defn}\ollabel{defn:normal-modal-logic}
یک منطق موجهات~$\Log L$ \emph{نرمال} است اگر و تنها اگر:
\begin{enumerate}
\item هر دو فرمول زیر را در بر داشته باشد:
\begin{align}
\tag{K} \Box(p \lif q) & \lif (\Box p \lif \Box q) \\
\tag{Dual} \Diamond p & \liff \lnot \Box \lnot p
\end{align}
\item و تحت ضرورت‌بخشی بسته باشد؛ یعنی هرگاه $!A \in \Log L$، آنگاه
  $\Box !A \in \Log L$ نیز.
\end{enumerate}
\end{defn}

با دو راه اصلی برای تعریف منطق‌های موجهات روبه‌رو خواهیم شد. یکی راه
نظریه‌اثباتی است: مجموعهٔ !!{formula}هایی که می‌توان در دستگاه‌های
!!{derivation} معینی استنتاج کرد. دیگری راه نظریه‌مدلی است: مجموعه‌های
!!{formula}هایی که در رده‌های معینی از مدل‌ها معتبرند. این وضع همانند
منطق گزاره‌ای معمولی و نیز منطق مرتبهٔ اول است. آنجا نیز می‌توان
مجموعه‌ای از !!{formula}ها را به دو شیوه مشخص کرد؛ برای نمونه، به‌صورت
مجموعهٔ همان‌گویی‌هایی که در همهٔ انتساب‌های ارزش صادق‌اند و به‌صورت
مجموعهٔ همهٔ موارد !!{derivable} در میان !!{formula}های یک دستگاه
!!{derivation}. برای منطق‌های موجهاتی نرمال، این شیوهٔ دوگانهٔ تعیین
منطق با استفاده از مدل‌های رابطه‌ای و دستگاه‌های اثبات اصل‌موضوعی از
نوع هیلبرتی ممکن است.
\end{document}
